\section{Boundary balance and local transfer identities}\label{j:local}

We first prove the local identities used by both cylinder observables. The boundary balance also shows that the total half-plane arch weight is finite; this finiteness justifies the vacuum limits in Section~\ref{c7:finite}. The sharper arch exponent is proved in Section~\ref{c7:boundary}.

We keep the lattice, ports and first-exit convention of the introduction.
In particular a walk traversing $N$ triangle centers has weight $x^N$,
where $x=(2+\sqrt2)^{-1/2}$. The local turning balance is the
Duminil-Copin--Smirnov parafermionic cancellation. The extension to
rhombic tiles uses the Yang--Baxter weights discussed by Glazman and
Manolescu \cite{DuminilCopinSmirnov2012,GlazmanManolescu2017}.

The boundary estimate ultimately obtained in Section~\ref{c7:boundary} is: in the half-plane bounded by a horizontal tessellation line, the total weight $W_R$ of boundary-to-boundary walks from a fixed port reaching Euclidean distance $R$ or more from that port satisfies
\[
W_R=R^{-1/4+o(1)}.
\]

\label{j:convex-balance}
Write $s=3/8$, $\alpha=\pi/8$, so $2x\cos(s\pi/3)=1$. For any finite convex tessellation polygon $D$ and starting port $a$, the sum over all SAWs to ports, with the extra factor $e^{isT}$ for total turning $T$, is 1. In fact, extend the nonbacktracking walk from vertex to vertex, with factors $xe^{\pm is\pi/3}$ summing to 1, and stop before traversing an already visited vertex. Ends of this tree due to first collision cancel: after a common incoming stem there are two ways around the newly formed simple loop, and the stem approaches the loop from outside. The turn at closure would be $\pm\pi/3$ (unused third branch outside means a convex corner); thus other loop vertices contribute $\pm5\pi/3$ and the turn from the stem onto the loop $\mp\pi/3$. The extra factor therefore sums to $2\cos(4\pi s/3)=0$. For the walks to ports $b$, $T=\Delta-\pi$ where $\Delta$ is the counterclockwise boundary tangent turn accumulated from $a$ to $b$ (close the walk along the boundary from $b$ to $a$, giving positive orientation and two extra quarter turns). In particular $\cos(sT)\ge\sin\alpha>0$. This turning calculation is the local parafermionic balance in path form. \cite{DuminilCopinSmirnov2012}

Let $H(d)$ denote the half-plane arch weight from a port to a port $d$ steps to the right. Reflection gives the same weight to the left; by convex exhaustions $2\sin\alpha\sum_{d>0}H(d)\le1$.

\subsection{An auxiliary cylinder and spectral equation}
Take a cylinder obtained by identifying modulo $M$ steps horizontally, tiled by rows of $M$ rhombi each, with basis $1,e^{i\theta}$ and potentially a different $\theta$ per row. Connections inside a tile link pairs of bond midpoints on distinct sides, using each at most once, noncrossingly. The empty weight is 1, a straight connection between opposite sides weighs $v$, the single turns of absolute angles $\theta,\pi-\theta$ weigh $u,d$ respectively (port directions normal to sides); and two simultaneous turns weigh $a$ for the pairing left--down and up--right, or $b$ for left--up and right--down. With $\phi=3\theta/8$, $s_k=\sin(\phi+k\alpha)$, $h=\sin(2\alpha)$, use the trigonometric SAW weights
\[
(u,d,v,a,b)=\frac{(h s_5,h s_0,s_0s_5,s_5s_6,-s_0s_7)}{s_2s_3}.
\]
At $\phi=\alpha$ these are $(x,x^2,x^2,x^2,0)$, exactly the original lattice rules (split each tile into the two equilateral triangles). The relation between rhombic geometry, parafermionic cancellation and
integrable loop weights was developed by Ikhlef and Cardy
\cite{IkhlefCardy2009}. Glazman's corrected general loop-weight formulas
\cite[Section~5]{Glazman2015} give a useful comparison; the five weights
and their critical specialization here are checked directly below.
See also the boundary two-point results in \cite{GlazmanManolescu2017}.

For a finite stack, close both rims (require unused ports). The partition $Z$ of \textbf{one noncontractible polygon only} (unoriented, unrooted) is the linear coefficient in fugacity $n$ of the full loop sum giving weight zero to contractible loops and weight $n$ per noncontractible loop. For real angles $0<\theta<\pi$, orient loops with factor $\exp(-iT/4)$ and twist $p$ to the power signed number of seam crossings (east positive); summing orientations effects precisely these loop weights for $n=p+p^{-1}$. Indeed disjoint arcs can be drawn in tiles with the indicated tangent turns; an embedded loop in the flat cylinder is either contractible with turn $\pm2\pi$ or has winding $\pm1$ and turn zero. For instance the latter follows by mapping conformally by an exponential to a round annulus and using the plane turning formula there. These identities then extend algebraically. 

The oriented matrix, with spins $0,\pm1$ on horizontal-flow and descending-flow bonds (respectively across the sides parallel to $e^{i\theta}$ and to $1$), has row indices (left, up), column indices (right, down), positive signs meaning left to right or up to down. Remove the phase $e^{i\theta(\text{left}-\text{right})/4}$ (telescoping). The resulting $R(\phi)$ conserves the sum, with blocks, $\lambda=e^{-2i\alpha}$:
\[
\begin{array}{c|c}
(s',s')\quad(s'=\pm1)& a\\
(s',0),(0,s')&\begin{pmatrix}v&u\\u&v\end{pmatrix}\\
(1,-1),(0,0),(-1,1)&
\begin{pmatrix}
 b&d\lambda& b\lambda^2+a\\
 d\lambda^{-1}&1&d\lambda\\
 b\lambda^{-2}+a&d\lambda^{-1}&b
\end{pmatrix}
\end{array}
\]
(the turns connecting left to down have magnitude $\theta$, those connecting left to up have magnitude $\pi-\theta$).

\subsection{Local exchange and fusion}

The matrix identities below have two uses. Exchange lets us interchange
two rows of local matrices. Fusion identifies the action on a smaller auxiliary
space; at a matching site the action on the quotient vanishes. We treat
these as meromorphic identities, including at the rank-deficient fusion
matrices.

Let $V=\mathbb C^3$ with basis $|0\rangle,|+\rangle,|-\rangle$, and use
subscripts to indicate an ordered pair of tensor factors. On a copy
$V_i$, let $s_i|k\rangle_i=k|k\rangle_i$. Write $P_{ij}$ for the
interchange of factors $i,j$, and put
\[
\Delta_R(\psi)=\sin(\psi+2\alpha)\sin(\psi+3\alpha).
\]
\begin{lemma}[Local exchange and fusion]\label{j:local-exchange-fusion}
The entries of $R(\psi)$ are meromorphic, with possible poles at
$\Delta_R(\psi)=0$. The following are identities of meromorphic matrix
functions:
\[
\begin{split}
R_{12}(t)R_{13}(t+\phi)R_{23}(\phi)
 &=R_{23}(\phi)R_{13}(t+\phi)R_{12}(t),\\
R_{12}(\phi)R_{21}(-\phi)&=I.
\end{split}
\]
In particular the second equality is a pointwise inverse identity
where both factors are regular, and $R_{12}(0)=P_{12}$.

For $\delta=2,3$, let $E_2=V$, $E_3=\mathbb C$, and define injective
maps $J_\delta:E_\delta\to V\otimes V$ by
\[
\begin{aligned}
J_2|0\rangle&=|00\rangle+x\sum_{k=\pm1}\lambda^k|k,-k\rangle,
&J_2|k\rangle&=x(|k,0\rangle+|0,k\rangle)\quad(k=\pm1),\\
J_3(1)&=\sum_{k=0,\pm1}\lambda^k|k,-k\rangle.
\end{aligned}
\]
On $E_2$ let $s_{\rm f}|k\rangle=k|k\rangle$, and on $E_3$ let
$s_{\rm f}=0$.
Then $R_{12}(\delta\alpha)=J_\delta J_\delta^*$ and, writing $J=J_\delta$,
\[
R_{23}(\phi)R_{13}(\phi+\delta\alpha)(J\otimes I)
 =(J\otimes I)R_{\rm f}(\phi),\qquad
R_{\rm f}(\phi)=
\begin{cases}
R(\phi+\alpha),&\delta=2,\\
I,&\delta=3.
\end{cases}
\]
The intertwining equality is again meromorphic in $\phi$.
\end{lemma}
\begin{proof}
At $0,2\alpha,3\alpha$ the five weights are, respectively,
$(1,0,0,1,0)$, $(x^2,x,x^2,0,x^2)$ and $(0,1,0,0,1)$.
The blocks give $R_{12}(0)=P_{12}$ and the two factorizations directly. For example,
the charge-zero column of $J_2$ is
$(x\lambda,1,x\lambda^{-1})^T$, and the charge-$\pm1$ columns are
$x(1,1)^T$. These are embeddings rather than isometries:
$J_2^*J_2=\operatorname{diag}(3-\sqrt2,2-\sqrt2,2-\sqrt2)$ in the
displayed basis, and $J_3^*J_3=3$.

For regular $\phi$ and $-\phi$, unitarity follows by block multiplication,
using $\lambda^4=-1$. With primes denoting $-\phi$, the required scalar
relations are
\[
aa'=1,\quad ab'+ba'+dd'=0,\quad ad'+d=a'd+d'=0,
 \quad vv'+uu'=1,\quad vu'+uv'=0
\]
by $s'_j=s_{8-j}$ and the sine products above. For instance the first
diagonal entry of the charge-zero product is
$aa'+\lambda^2(ab'+ba'+dd')$, and its adjacent entry is
$\lambda(ad'+d)$. Since these are rational trigonometric identities,
unitarity follows meromorphically as stated.

Writing $B(i,l;j,k;\phi)$ for entries of $R$, the same blocks give
\[
B(i,l;r,k;\phi+3\alpha)=\lambda^{k-l}B(l,-r;k,-i;-\phi)
\]
(weights swap $u,d$ and $a,b$, keeping $v$). Thus for the single column intertwining, the coefficient is
\[
\sum_r B(j,m;-r,l;\phi) B(i,l;r,k;\phi+3\alpha)\lambda^r
=\lambda^i\delta_{j,-i}\delta_{m,k},\qquad l=r+k-i,
\]
with out-of-range entries zero.

For the other intertwining, denote the five weights at $\phi+2\alpha$ by $U,D,V,A,B_0$, at $\phi+\alpha$ by $w,X,g,e,f$. To detail the block multiplication, the components with inputs fused spin $r$ and site spin $k$, outputs $i,j,m$ ($i+j+m=r+k$) are
\[
\sum_t J_{t,r-t;r}\, R(\phi+2\alpha)_{i,t+k-i;t,k}\,R(\phi)_{j,m;r-t,t+k-i}
 = J_{i,j;i+j} R(\phi+\alpha)_{i+j,m;r,k}.
\]
It suffices (opposite spins conjugate the phases) to check the following scalar identities for each group after substitution and $\lambda^2+\lambda^{-2}=0$:
{\small
\[
\begin{array}{c|l}
(r,k) &\\ \hline
(1,1) & Va=Av+Uu=e,\quad Au+Uv=0\\
(1,0)& Da=w,\quad x(Uv+u)=w,\quad Uu+v=Dd+V=g,\quad Db+Vd=0\\
(0,1)& B_0 a=g,\quad x D u+Vv=g,\quad x Dv+Vu=xw,\quad x(Aa+B_0 b)+Ud=0,\quad xB_0 d+U=xw\\
(1,-1)& B_0 v=Dd+V b=f,\quad B_0 u=x(D+Vd)=X,\quad Va=Av+Uu=e,\quad Au+Uv=0\\
(0,0)& x Ua+d=x X,\quad D u+x Vv=x,\quad Dv+x Vu=x X
\end{array}
\]
}
For clarity $U,D,V,A,B_0$ have numerators $(h s_7,h s_2,s_2 s_7,-s_0s_7,s_2s_1)$ over $s_4s_5$, and $w,X,g,e,f$ numerators $(h s_6,h s_1,s_1s_6,s_6s_7,s_1s_0)$ over $s_3s_4$. Direct common factors cancel; the remaining identities reduce to
\[
\begin{gathered}
 s_2s_6=h^2-s_0^2=x(s_0s_7+s_4s_5),
 \quad s_1s_5=h^2-s_7^2=x(s_0s_7+s_2s_3),\\
 s_1-s_7=s_0/x,\quad s_6-s_0=s_7/x,\\
 s_0s_1+s_4s_5=s_2s_3+s_6s_7=s_1s_2+s_5s_6=h^2/x,\\
 s_1s_6-s_0s_7=h^2 x,\quad s_3s_4-s_0s_7=h^2/x,\quad
 x(s_5s_6-s_0s_1)=s_3s_7,\quad x(s_1s_2-s_6s_7)=s_0s_4,
\end{gathered}
\]
all by product-to-sum.

For regular real parameters the matrices are Hermitian. The intertwining
and its adjoint therefore give the three-space exchange at
$t=2\alpha,3\alpha$; the swap gives it at $t=0$. Interchange spaces
2 and 3 in these equalities and apply unitarity at a generic $\phi$.
This gives exchange also at $t+\phi=0,2\alpha,3\alpha$.

It remains to see why these special values determine the identity.
Every entry of $\Delta_R(\psi)R(\psi)$ has Fourier exponents between
$-2$ and $2$. Moreover, if $Z|k\rangle=(-1)^{k^2}|k\rangle$, then
$R(\psi+\pi)=Z_1R(\psi)Z_1$: only the single-turn weights change
sign. After multiplying the exchange difference by
$\Delta_R(t)\Delta_R(t+\phi)\Delta_R(\phi)$, each entry consequently
has exponents between $-4$ and $4$ in $e^{it}$, all of one parity.
After a monomial factor it is a polynomial in $e^{2it}$ of degree at
most $4$. The six values
$0,2\alpha,3\alpha,-\phi,2\alpha-\phi,3\alpha-\phi$ are distinct
modulo $\pi$ for generic real $\phi$, so the polynomial vanishes.
Continuation in the two parameters proves the meromorphic identity.
\end{proof}

\begin{lemma}[Restriction and quotient at fusion]\label{j:local-restriction}
Fix $\delta\in\{2,3\}$, put $W_\delta=\operatorname{im}J_\delta
\subset V_1\otimes V_2$, and for another factor $V_k$ set
\[
F_{\delta,k}(\psi)=R_{2k}(\psi)R_{1k}(\psi+\delta\alpha).
\]
At regular parameters the subspace $W_\delta\otimes V_k$ is invariant,
and transporting the restriction through $J_\delta\otimes I$ gives $R_{\rm f}(\psi)$ from
Lemma~\ref{j:local-exchange-fusion}. At $\psi=0$ the full image of
$F_{\delta,k}(0)$ lies in this subspace. Hence the induced operator on
$\bigl((V_1\otimes V_2)/W_\delta\bigr)\otimes V_k$ is zero there and, near
$z=e^{2i\psi}=1$, is holomorphic and divisible by $z-1$.

The twist restricts by
$p^{s_1+s_2}J_2=J_2p^{s_{\rm f}}$ and
$p^{s_1+s_2}J_3=J_3$ for every $p\ne0$.
There is also the meromorphic site intertwining
\begin{equation}\label{j:site-fusion}
R_{a1}(\psi)R_{a2}(\psi+\delta\alpha)(I_a\otimes J_\delta)
 =(I_a\otimes J_\delta)
\begin{cases}
R_{a{\rm f}}(\psi+\alpha),&\delta=2,\\
I_a,&\delta=3.
\end{cases}
\end{equation}
\end{lemma}
\begin{proof}
Invariance and the restricted action are the preceding intertwining.
Write $J_{\delta,1k}$ for the embedding on factors 1 and $k$, with
the other factor unchanged.
At the matching point the factorization supplies the stronger inclusion
\[
F_{\delta,k}(0)=P_{2k}R_{1k}(\delta\alpha)
 =P_{2k}J_{\delta,1k}J_{\delta,1k}^*,
\qquad
\operatorname{im}F_{\delta,k}(0)\subseteq W_\delta\otimes V_k.
\]
Indeed $P_{2k}$ carries the image of $J_{\delta,1k}$ to that of
$J_{\delta,12}$. The quotient operator is holomorphic near zero because
the local parameters $0,\delta\alpha$ are regular. Its vanishing gives
the claimed factor in the local coordinate $z-1$. The twist formulas
follow because every term of $J_2|s_{\rm f}\rangle$ has total spin
$s_{\rm f}$ and every term of $J_3$ has total spin zero.

For real regular parameters the blocks give
$\overline{R_{ij}}=R_{ji}$ and $P_{12}\overline{J_\delta}=J_\delta$.
Conjugating the row intertwining and exchanging spaces 1 and 2 gives
\eqref{j:site-fusion}; continuation gives its meromorphic form.
\end{proof}


\subsection{A Fourier identity for the scalar kernels}

The finite vacuum normalization and the later boundary moment problem
use the same elementary transform. Recording it here makes both
calculations independent of their order.

\begin{lemma}\label{j:fourier-kernel}
For $0<\beta<\pi$ and $|\Im y|<\pi-\beta$,
\[
 \int_{\mathbb R}e^{iky}\frac{\sinh(\beta k)}{\sinh(\pi k)}\,dk
 =\frac{\sin\beta}{\cosh y+\cos\beta}.
\]
The integral and its parameter derivatives converge locally uniformly
on the indicated strip.
\end{lemma}
\begin{proof}
For real $y>0$, close the contour by rectangles in the upper half-plane
with upper edges at half-integer heights. The ratio of hyperbolic sines
decays exponentially on the vertical edges as their real coordinates
tend to infinity; the factor $e^{iky}$ controls the upper edge as its
height tends to infinity. The residues at $k=in$, $n\ge1$, give
\[
 2\sum_{n\ge1}(-1)^{n+1}\sin(n\beta)e^{-ny}
 =\frac{\sin\beta}{\cosh y+\cos\beta}.
\]
The geometric series proves the last equality. Evenness extends it to
negative real $y$, and dominated convergence covers $y=0$.
Exponential decay in $|k|$ proves local uniform convergence and
parameter differentiation on smaller strips. Analytic continuation
then gives the asserted domain.
\end{proof}
